% ============================================================================
% sec12.tex — §12 Thermal Radiation, Blackbody Laws (Planck, Wien,
%             Stefan-Boltzmann), View Factors, Radiation Shields & Selective Surfaces
% Target: 8 pages | Figures: fig:12:planck-spectra, fig:12:shields-viewfactor
% ============================================================================
\binderchapter{12}{Thermal Radiation, Blackbody Laws, View Factors, Radiation Shields \& Two-Band Selective Surfaces}{Thermal Radiation}{sec:12}
\scopebanner{TIER-DEPENDENT}[2027 Rule 5.c.iii (All Levels Core) \& Rule 5.c.vi (States/Nationals: Radiant Exitance, Blackbody Radiation, Stefan--Boltzmann Law) \textbar\ Sources: \src{S01}, \src{S07}, \src{S12}, \src{S19}, \src{S25}]

\idxtopic{radiation!Planck blackbody distribution}
\idxtopic{radiation!Wien's displacement law}
\idxtopic{radiation!Stefan-Boltzmann law and radiant exitance}
\idxtopic{radiation!Lambertian radiance $L=\sigma T^4/\pi$}
\idxtopic{view factor!reciprocity and coaxial disks}
\idxtopic{radiation shield!multi-layer foil suppression}
\idxtopic{selective surfaces!solar/NIR absorptance vs MIR emittance}

% ----------------------------------------------------------------------------
% L1: FORMULA & DECISION CARD
% ----------------------------------------------------------------------------
\begin{formulacard}[§12 Master Blackbody Radiation, View-Factor, Shield \& Selective-Surface Decision Card \src{S07,S12,S25}]
\begin{ftab}
$\alpha + \rho + \tau = 1, \quad \alpha_\lambda + \rho_\lambda + \tau_\lambda = 1$
  & $\alpha$: absorptance; $\rho$: reflectance; $\tau$: transmittance (all $\in [0, 1]$); for an \textbf{opaque} surface ($\tau=0$): $\alpha + \rho = 1$
  & \flag{OPTICAL}
  & Energy conservation for incident irradiance $G$ $(\si{\watt\per\square\meter})$ at any surface \idxformula{radiation balance!$\alpha+\rho+\tau=1$} \\
$M_{\lambda,b}(\lambda, T) = \dfrac{c_1}{\lambda^5 \left[\exp\!\left(\dfrac{c_2}{\lambda T}\right) - 1\right]}$
  & $M_{\lambda,b}$ (or $e_{b\lambda}$): spectral blackbody exitance $(\si{\watt\per\square\meter\per\micro\meter})$; $c_1 = 2\pi h c_0^2 = \SI{3.74177e8}{\watt\micro\meter\tothe{4}\per\square\meter}$; $c_2 = \dfrac{h c_0}{k_B} = \SI{14387.77}{\micro\meter\kelvin}$
  & \TIERf
  & Planck distribution (1900); $T$ MUST be in $\si{\kelvin}$, $\lambda$ in $\si{\micro\meter}$ \idxformula{Planck's law!$M_{\lambda,b}(\lambda,T)$} \\
$\lambda_{\max} T = b_{\mathrm{Wien}} = \SI{2897.772}{\micro\meter\kelvin}$
  & $\lambda_{\max}$: peak wavelength of $M_{\lambda,b}$ $(\si{\micro\meter})$; $b_{\mathrm{Wien}} = c_2 / 4.965114$ $(\si{\micro\meter\kelvin})$ [frequency peak: $\nu_{\max}/T = \SI{5.8789e10}{\hertz\per\kelvin}$]
  & \TIERf
  & Sun ($\SI{5778}{\kelvin}$): $\SI{0.502}{\micro\meter}$; 125 W Lamp ($\SI{2400}{\kelvin}$): $\SI{1.207}{\micro\meter}$; Water ($\SI{300}{\kelvin}$): $\SI{9.66}{\micro\meter}$ \idxformula{Wien's displacement law!$\lambda_{\max}T=b$} \\
$M_b(T) = \sigma T^4, \quad M(T) = \varepsilon \sigma T^4$
  & $M$ (or $e$): hemispherical radiant exitance $(\si{\watt\per\square\meter})$; $\sigma = \dfrac{2\pi^5 k_B^4}{15 c_0^2 h^3} = \SI{5.670374e-8}{\watt\per\square\meter\per\kelvin\tothe{4}}$; $\varepsilon \in [0, 1]$: total emittance
  & \TIERf
  & Stefan--Boltzmann law ($\dot{Q}_{\mathrm{emit}} = \varepsilon\sigma A T^4$); $T$ MUST be in absolute $\si{\kelvin}$! \idxformula{Stefan-Boltzmann law!$M=\varepsilon\sigma T^4$} \\
$L = I = \dfrac{M}{\pi} = \dfrac{\varepsilon \sigma T^4}{\pi}$
  & $L$ (or $I$): directional radiance / intensity $(\si{\watt\per\square\meter\per\steradian})$ of a diffuse (Lambertian) emitter; factor is $\pi$ (NOT $2\pi$!)
  & \TIERf
  & 2027 Cornell Invitational \src{S25} \& Lienhard Eq.~(10.7b) p.~551 \src{S12} \idxformula{Lambertian radiance!$L=\varepsilon\sigma T^4/\pi$} \\
$\varepsilon_\lambda(T,\theta,\phi) = \alpha_\lambda(T,\theta,\phi)$
  & Exact Kirchhoff's law: holds at identical wavelength $\lambda$; total $\alpha_{\mathrm{lamp}} \ne \varepsilon_{\mathrm{MIR}}$ when $T_{\mathrm{lamp}} \approx \SI{2400}{\kelvin} \gg T_{\mathrm{surf}} \approx \SI{310}{\kelvin}$!
  & \flag{SELECTIVE}
  & Two-band selective absorbers ($\alpha_{\mathrm{NIR}}/\varepsilon_{\mathrm{MIR}} \gg 1$) \& greenhouse glazings \idxformula{Kirchhoff's law!$\varepsilon_\lambda=\alpha_\lambda$} \\
$\sum_{j=1}^N F_{i\to j} = 1, \quad A_i F_{i\to j} = A_j F_{j\to i}$
  & $F_{i\to j}$: geometric view factor (fraction of diffuse radiation leaving $i$ that strikes $j$ directly); convex/flat surface has $F_{i\to i} = 0$
  & \flag{GEOMETRY}
  & Enclosure view-factor algebra \& coaxial lamp-to-aperture interception \idxformula{view factor!summation and reciprocity} \\
$\dot{Q}_{1\to 2} = A_1 \mathcal{F}_{1\to 2}\,\sigma\left(T_1^4 - T_2^4\right)$
  & $\mathcal{F}_{1\to 2} = \varepsilon_1$ (small body $A_1 \ll A_2$); $\mathcal{F}_{1\to 2} = \left[\dfrac{1}{\varepsilon_1} + \dfrac{A_1}{A_2}\!\left(\dfrac{1}{\varepsilon_2}-1\right)\right]^{-1}$ (enclosed body, $F_{12}=1$)
  & \GRAY
  & Two-body diffuse-gray exchange (Lienhard Eqs.~10.25--10.30 \src{S12}) \idxformula{gray-body radiation exchange!$\dot{Q}_{12}=A_1\mathcal{F}_{12}\sigma(T_1^4-T_2^4)$} \\
$\dot{Q}_{N\text{ shields}} = \dfrac{\sigma A\left(T_1^4 - T_2^4\right)}{(N+1)\left(\dfrac{2}{\varepsilon} - 1\right)} = \dfrac{\dot{Q}_0}{N+1}$
  & $N$: number of parallel radiation shields between plates when all surfaces have emittance $\varepsilon$; for unequal $\varepsilon_k$, resistances add in series
  & \flag{SHIELDS}
  & Multi-layer aluminum-foil insulation ($\varepsilon=0.04$ cuts $\dot{Q}_{\mathrm{rad}}$ by $41.7(N+1)\times$!) \idxformula{radiation shields!$\dot{Q}_N=\dot{Q}_0/(N+1)$} \\
$h_{\mathrm{rad}} = 4\sigma T_m^3 \mathcal{F}_{1\to 2}, \quad R_{\mathrm{rad}} = \dfrac{1}{h_{\mathrm{rad}}A_1}$
  & $T_m = \frac{1}{2}(T_1 + T_2)$ $(\si{\kelvin})$; valid within $<1\%$ when $\frac{1}{4}(\Delta T/T_m)^2 \ll 1$; at $T_m=\SI{300}{\kelvin}$, $4\sigma T_m^3 = \SI{6.12}{\watt\per\square\meter\per\kelvin}$
  & \flag{LINEARIZED}
  & Parallel convection + radiation networks ($h_{\mathrm{eff}} = h_{\mathrm{conv}} + h_{\mathrm{rad}}$) \idxformula{linearized radiation coefficient!$h_{\mathrm{rad}}=4\sigma T_m^3\mathcal{F}_{12}$} \\
\end{ftab}
\end{formulacard}

% ----------------------------------------------------------------------------
% L2: COMPACT EXPLANATION & THREE-SPECTRUM SEPARATION
% ----------------------------------------------------------------------------
\section{Blackbody Radiation, Spectral Separation \& Two-Band Selectivity}
\label{sec:12.1}

\begin{tierbox}[2027 Rule \texttt{5.c.iii} (All Levels) vs.\ Rule \texttt{5.c.vi} (States \& Nationals Only) \src{S01}]
Qualitative radiation mechanisms, shiny low-emissivity barriers, and 2027 heat-lamp collection appear at \textbf{all competition levels} under Rule \texttt{5.c.iii} and Part~II device testing. Quantitative calculations of \textbf{radiant exitance ($M$), Planck/Wien blackbody spectra, and the Stefan--Boltzmann law ($M = \varepsilon\sigma T^4$)} are explicitly listed under Rule \texttt{5.c.vi} (\emph{State and National levels only}, though frequent at major invitationals such as the 2026--2027 Cornell Invitational \src{S25}).
\end{tierbox}

\subsection{Planck Distribution, Wien's Displacement Law, and Lambertian Radiance}
Every body above $\SI{0}{\kelvin}$ continuously emits electromagnetic thermal radiation (\SIrange{0.1}{1000}{\micro\meter}, Lienhard Table~1.2 p.~28 \src{S12}). An ideal \textbf{blackbody} (realized experimentally by a small aperture in an isothermal cavity or \emph{hohlraum}, where multiple internal reflections trap all incident rays) absorbs $100\%$ of incident radiation at all wavelengths and directions ($\alpha_\lambda = 1,\; \rho_\lambda = \tau_\lambda = 0$) and emits the maximum possible spectral flux allowed by thermodynamic equilibrium:
\begin{enumerate}
  \item \textbf{Wien's Displacement Law (1893) and the Wavelength--Frequency Peak Trap:} Differentiating Planck's distribution $M_{\lambda,b}(\lambda,T)$ with respect to $\lambda$ yields the transcendental equation $(5 - x)e^x = 5$ where $x = h c_0 / (\lambda_{\max} k_B T) = 4.96511423$, giving $\lambda_{\max} T = b_{\mathrm{Wien}} = \SI{2897.772}{\micro\meter\kelvin}$ \src{S07,S12}. As plotted in \fref{fig:12:planck-spectra}, this creates a crisp \textbf{three-spectrum wavelength separation} critical to the 2027 collector:
  \begin{itemize}
    \item \textbf{Solar Spectrum ($T_\odot \approx \SI{5772}{\kelvin}$--$\SI{5778}{\kelvin}$):} $\lambda_{\max} = \frac{2897.8}{5778} = \SI{0.5015}{\micro\meter}$ ($\SI{502}{\nano\meter}$, visible green-yellow; $\sim 98\%$ of solar energy lies at $\lambda < \SI{2.7}{\micro\meter}$).
    \item \textbf{2027 Incandescent 100--125 W Heat-Lamp Tungsten Filament ($T_{\mathrm{fil}} \approx \SI{2400}{\kelvin}$):} $\lambda_{\max} = \frac{2897.8}{2400} = \mathbf{\SI{1.2074}{\micro\meter}}$ (\textbf{near-infrared}; for $T_{\mathrm{fil}} \in [\SI{2200}{\kelvin}, \SI{2500}{\kelvin}]$, $\lambda_{\max} \in [\SI{1.159}{\micro\meter}, \SI{1.317}{\micro\meter}]$). Exactly $25.0\%$ of total blackbody emission always lies at $\lambda < \lambda_{\max}$, and $75.0\%$ lies at $\lambda > \lambda_{\max}$.
    \item \textbf{Room-Temperature Water \& Collector Walls ($T_w \approx \SI{300}{\kelvin}$--$\SI{318}{\kelvin}$):} $\lambda_{\max} = \frac{2897.8}{300} = \mathbf{\SI{9.659}{\micro\meter}}$ ($\SI{9.11}{\micro\meter}$ at $\SI{45}{\celsius}$, \textbf{mid/thermal infrared}, spanning $\SIrange{4}{40}{\micro\meter}$).
  \end{itemize}
  \emph{Caution on Frequency Peak $\nu_{\max}$:} Because $M_\nu = M_\lambda |d\lambda/d\nu| = (\lambda^2/c_0)M_\lambda$, the peak of Planck's distribution per unit frequency $M_{\nu,b}(\nu,T)$ occurs at $x_\nu = h\nu_{\max}/(k_B T) = 2.821439$, so $\nu_{\max}/T = \SI{5.8789e10}{\hertz\per\kelvin}$ and $\lambda_{\max}\nu_{\max} = \frac{2.8214}{4.9651}c_0 = 0.56825\,c_0 \ne c_0$!
  \item \textbf{Lambert's Cosine Law and Directional Radiance ($L = M/\pi$):} A blackbody or ideal diffuse surface emits with uniform \textbf{directional radiance} (intensity) $L = I$ $[\si{\watt\per\square\meter\per\steradian}]$ in all directions $(\theta, \phi)$, so the heat flux emitted per unit solid angle $d\Omega = \sin\theta\,d\theta\,d\phi$ at polar angle $\theta$ from the surface normal is $dM = L \cos\theta\,d\Omega$ (because the projected area seen from angle $\theta$ is $dA\cos\theta$). Integrating over the upper hemisphere ($\phi \in [0, 2\pi],\; \theta \in [0, \pi/2]$) gives $\int_0^{2\pi}d\phi \int_0^{\pi/2}\cos\theta\sin\theta\,d\theta = 2\pi\left(\frac{1}{2}\right) = \pi\ \si{\steradian}$ (Lienhard Eq.~10.7b, p.~551 \src{S12}):
  \begin{equation}
    M(T) = \pi L(T) = \varepsilon \sigma T^4 \implies L(T) = \frac{\varepsilon \sigma T^4}{\pi} \quad \left[\si{\watt\per\square\meter\per\steradian}\right]
    \label{eq:12:lambert-radiance}
  \end{equation}
\end{enumerate}

% ---------------- FIGURE 32: fig:12:planck-spectra ---------------------------
\begin{center}
\begin{tikzpicture}[font=\scriptsize]
  \begin{axis}[
    width=0.92\linewidth, height=4.8cm,
    xmode=log, ymode=log,
    xlabel={Wavelength $\lambda$ $(\si{\micro\meter})$ — Showing Short-Wave Lamp Band ($\lambda < \SI{2.7}{\micro\meter}$) vs.\ Long-Wave Thermal IR Band ($\lambda > \SI{4}{\micro\meter}$)},
    ylabel={Spectral Exitance $M_{\lambda,b}(\lambda,T)$ $(\si{\watt\per\square\meter\per\micro\meter})$},
    xmin=0.2, xmax=35, ymin=1, ymax=1e8,
    xtick={0.2, 0.4, 0.7, 1.21, 2.7, 5, 9.66, 20, 35},
    xticklabels={$0.2$, $0.4$, $0.7$, $1.21$, $2.7$, $5.0$, $9.66$, $20$, $35$},
    ytick={1, 1e2, 1e4, 1e6, 1e8},
    grid=both, major grid style={gray!30}, minor grid style={gray!12},
    legend pos=north east, legend style={font=\scriptsize, fill=white, draw=navy},
    clip=true
  ]
    % Shaded glass/BoPET cutoff at 2.7 um
    \fill[teal!10] (axis cs:0.35,1) rectangle (axis cs:2.7,1e8);
    \draw[teal, very thick, dashed] (axis cs:2.7,1) -- (axis cs:2.7,1e8);
    \node[teal, font=\bfseries\scriptsize, align=center] at (axis cs:0.85, 2.8e7) {Glass / BoPET Window\\[-1pt]Transmits ($\tau_{\mathrm{NIR}} \approx 0.90$)};
    \node[rust, font=\bfseries\scriptsize, align=center] at (axis cs:6.8, 2.8e7) {Glass / BoPET Opaque\\[-1pt]Traps Mid-IR ($\tau_{\mathrm{MIR}} < 0.20$)};

    % Planck curve T = 5778 K (Sun at surface)
    \addplot[amber, very thick, domain=0.2:30, samples=36] {3.74177e8 / (x^5 * (exp(14387.77/(x*5778)) - 1))};
    \addlegendentry{Solar Blackbody ($T=\SI{5778}{\kelvin},\; \lambda_{\max}=\SI{0.502}{\micro\meter}$)}

    % Planck curve T = 2400 K (2027 125 W Heat Lamp Tungsten Filament)
    \addplot[rust, ultra thick, domain=0.3:35, samples=36] {3.74177e8 / (x^5 * (exp(14387.77/(x*2400)) - 1))};
    \addlegendentry{2027 Heat Lamp ($T=\SI{2400}{\kelvin},\; \lambda_{\max}=\SI{1.207}{\micro\meter}$)}

    % Planck curve T = 300 K (Room Water)
    \addplot[navy, ultra thick, domain=2.0:35, samples=36] {3.74177e8 / (x^5 * (exp(14387.77/(x*300)) - 1))};
    \addlegendentry{Cup Water ($T=\SI{300}{\kelvin},\; \lambda_{\max}=\SI{9.659}{\micro\meter}$)}

    % Wien displacement locus: M_lambda_max = 1.2865e-11 * T^5 = 1.2865e-11 * (2897.77/x)^5 = 2.625e6 / x^5
    \addplot[black, thick, densely dotted, domain=0.45:12, samples=40] {2.625e6 / (x^5)};
    \addlegendentry{Wien Peak Locus ($\lambda_{\max} = \SI{2897.8}{\micro\meter\kelvin}/T$)}

    \fill[amber] (axis cs:0.5015, 8.27e7) circle (2pt);
    \fill[rust]  (axis cs:1.2074, 1.024e6) circle (2.2pt) node[above right, text=rust, font=\bfseries\scriptsize] {Lamp Peak $\SI{1.21}{\micro\meter}$};
    \fill[navy]  (axis cs:9.659,  31.2) circle (2.2pt) node[above right, text=navy, font=\bfseries\scriptsize] {Water Peak $\SI{9.66}{\micro\meter}$};
  \end{axis}
\end{tikzpicture}
\captionof{figure}{Planck blackbody spectral exitance $M_{\lambda,b}(\lambda,T)$ for the Sun ($\SI{5778}{\kelvin}$), the 2027 125~W tungsten heat lamp ($\SI{2400}{\kelvin}$, peaking at $\lambda_{\max} = \SI{1.207}{\micro\meter}$ in the near-IR), and room-temperature water ($\SI{300}{\kelvin}$, peaking at $\SI{9.66}{\micro\meter}$ in the mid-IR), overlaid with the $\SI{2.7}{\micro\meter}$ selective transmission cutoff of glass and BoPET glazing (Lienhard §10.1 \& §10.6 \src{S12}).}
\label{fig:12:planck-spectra}
\idxfig{Planck blackbody spectrum!solar vs 2400 K heat lamp vs 300 K water}
\end{center}

\subsection{Kirchhoff's Law and Two-Band Wavelength-Selective Surfaces (Lienhard Tables 10.1 \& 10.5)}
Kirchhoff's law states that at any given wavelength $\lambda$ and direction $(\theta,\phi)$, spectral directional emittance equals spectral directional absorptance: $\varepsilon_\lambda(T,\theta,\phi) = \alpha_\lambda(T,\theta,\phi)$ (Lienhard Eq.~10.8a, p.~551 \src{S12}). However, when an object at $T_{\mathrm{surf}} \approx \SI{310}{\kelvin}$ is irradiated by a high-temperature source at $T_{\mathrm{source}} \approx \SI{2400}{\kelvin}$ (or $\SI{5778}{\kelvin}$), the absorbed radiation is governed by the surface's short-wave near-IR absorptance $\alpha_{\mathrm{NIR}} \approx \varepsilon_\lambda(\lambda \approx \SI{1.2}{\micro\meter})$, whereas its emitted thermal loss is governed by its long-wave mid-IR emittance $\varepsilon_{\mathrm{MIR}} \approx \varepsilon_\lambda(\lambda \approx \SI{9.5}{\micro\meter})$:
\begin{equation}
  q''_{\mathrm{net,absorbed}} = \alpha_{\mathrm{NIR}}\,G_{\mathrm{lamp}} - \varepsilon_{\mathrm{MIR}}\,\sigma\left(T_{\mathrm{surf}}^4 - T_{\mathrm{amb}}^4\right) - h_{\mathrm{conv}}\left(T_{\mathrm{surf}} - T_{\mathrm{amb}}\right)
  \label{eq:12:selective-balance}
\end{equation}
\tref{tab:12:selective-emissivity} summarizes verified short-wave absorptances $\alpha_{\mathrm{SW/NIR}}$ and long-wave thermal emittances $\varepsilon_{\mathrm{MIR}}$ from Lienhard Tables~10.1 (p.~546) and 10.5 (p.~598) \src{S12}:

\begin{table}[htbp]
\centering
\caption{Verified Short-Wave (Solar/Near-IR) Absorptance $\alpha_{\mathrm{SW/NIR}}$, Reflectance $\rho_{\mathrm{SW/NIR}}$, and Long-Wave Mid-IR Emittance $\varepsilon_{\mathrm{MIR}}$ near $\SIrange{300}{313}{\kelvin}$ (Lienhard Tables 10.1 p.~546 \& 10.5 p.~598 \src{S12}).}
\label{tab:12:selective-emissivity}
\footnotesize
\setlength{\tabcolsep}{3.0pt}
\renewcommand{\arraystretch}{1.12}
\begin{tabularx}{\linewidth}{@{}L{0.34\linewidth} r r r r X@{}}
\toprule
\textbf{Surface / Coating (Lienhard Tables 10.1 \& 10.5)} & \textbf{$\alpha_{\mathrm{SW/NIR}}$} & \textbf{$\rho_{\mathrm{SW/NIR}}$} & \textbf{$\varepsilon_{\mathrm{MIR}}$} & \textbf{$\alpha_{\mathrm{SW}}/\varepsilon_{\mathrm{MIR}}$} & \textbf{Radiative Classification \& 2027 Role} \\
\midrule
\textbf{Black Chrome on Ni plate} (Table 10.5) & $0.95$ & $0.05$ & $0.09$ & $\mathbf{10.56}$ & Selective absorber: absorbs lamp, blocks MIR loss \\
\textbf{Sputtered cermet on steel} (Table 10.5) & $0.96$ & $0.04$ & $0.16$ & $\mathbf{6.00}$ & High-efficiency selective solar/NIR absorber \\
\textbf{$\mathrm{CuO}$ on copper (Ebanol C)} (Table 10.5) & $0.90$ & $0.10$ & $0.16$ & $\mathbf{5.63}$ & Selective black-oxide copper foil collar \\
\textbf{Nickel black on steel} (Table 10.5) & $0.81$ & $0.19$ & $0.17$ & $\mathbf{4.76}$ & Selective metallic absorber \\
\textbf{Copper, polished electrolytic} (Tables 10.1, 10.5) & $0.30$ & $0.70$ & $0.02\text{--}0.04$ & $7.50$ & Weak solar absorber, ultra-low MIR emitter \\
\textbf{Aluminum foil, bright polished} (Tables 10.1, 10.5) & $0.04\text{--}0.09$ & $0.91\text{--}0.96$ & $0.03\text{--}0.04$ & $1.0\text{--}2.2$ & Broadband mirror: ideal cone \& radiation shield \\
\textbf{Aluminum sheet, heavily oxidized} (Table 10.1) & $0.25$ & $0.75$ & $0.20\text{--}0.33$ & $\sim 1.0$ & Oxide layer degrades both reflectance \& shield $\varepsilon$ \\
\textbf{Carbon black in acrylic / Matte black paint} & $0.94$ & $0.06$ & $0.83\text{--}0.95$ & $1.0\text{--}1.13$ & Non-selective blackbody: needs glazing cover \\
\textbf{Carbon lampsoot / Candle soot} (Table 10.1) & $0.96$ & $0.04$ & $0.95$ & $1.01$ & Reference broadband blackbody ($\varepsilon = 0.95$) \\
\textbf{Liquid water layer (thickness $\ge\SI{0.1}{\milli\meter}$, \SI{40}{\celsius})} & $0.78\text{--}0.85^*$ & $0.02$ & $\mathbf{0.96}$ & $0.85$ & Direct NIR volume absorption; black in $8\text{--}14\,\si{\micro\meter}$! \\
\textbf{Smooth glass / Pyrex} (Table 10.1, §10.6 p.\ 599) & $\tau_{\mathrm{NIR}}=0.91$ & $0.08$ & $\mathbf{0.94}$ & Greenhouse & Transmits $\lambda<\SI{2.7}{\micro\meter}$, opaque in mid-IR ($\tau_{\mathrm{MIR}}<0.02$) \\
\textbf{White acrylic paint / White paper} (Tables 10.1, 10.5) & $0.26 / 0.22$ & $0.74 / 0.78$ & $\mathbf{0.90 / 0.96}$ & $\mathbf{0.23\text{--}0.29}$ & Selective \emph{radiative cooler}: reflects visible, emits MIR! \\
\textbf{Zinc oxide ($\mathrm{ZnO}$) white coating} (Table 10.5) & $0.15$ & $0.85$ & $\mathbf{0.93}$ & $\mathbf{0.16}$ & Stays cooler than air under lamp; never use in cone! \\
\bottomrule
\multicolumn{6}{@{}l@{}}{\scriptsize $^*$Planck-integrated absorptance of a $\SI{35}{\milli\meter}$ deep water column in a 9 oz cup under a $\SI{2400}{\kelvin}$ glass-filtered heat lamp (Hale \& Querry 1973 \src{S19}).}
\end{tabularx}
\end{table}

% ----------------------------------------------------------------------------
% L3: VIEW FACTORS, RADIATION SHIELDS & PLANETARY EQUILIBRIUM
% ----------------------------------------------------------------------------
\section{View-Factor Algebra, Gray-Body Networks \& Radiation Shields}
\label{sec:12.2}

\subsection{Coaxial-Disk View Factor ($F_{1\to 2}$) for the 2027 Heat Lamp and Collector Aperture}
For two coaxial parallel circular disks of radii $r_1$ (heat-lamp reflector face, $r_1 = \SI{6.35}{\centi\meter}$ for a $5.0\ \text{in}$ R40 bulb) and $r_2$ (collector top aperture $r_2 = \SI{15.0}{\centi\meter}$, or bare 9 oz cup mouth $r_{\mathrm{cup}} = \SI{3.75}{\centi\meter}$) separated by vertical distance $d = h_{\mathrm{bulb}} - z_{\mathrm{opening}}$, Lienhard Table~10.3 Case~3 (p.~562 \src{S12}) gives the exact analytical view factor:
\begin{equation}
  R_1 = \frac{r_1}{d}, \qquad R_2 = \frac{r_2}{d}, \qquad X = 1 + \frac{1 + R_2^2}{R_1^2}, \qquad F_{1\to 2} = \frac{1}{2}\!\left[X - \sqrt{X^2 - 4\left(\frac{R_2}{R_1}\right)^{\!2}}\right]
  \label{eq:12:coaxial-disks}
\end{equation}
Furthermore, by the enclosure summation rule (Lienhard Example~10.3, p.~565 \src{S12}), if disk $1$ (base radius $r_1$) and disk $2$ (top radius $r_2$) are connected by an enclosing conical frustum $s$, the fraction of diffuse radiation leaving disk $1$ that strikes the conical wall $s$ is simply $F_{1\to s} = 1 - F_{1\to 2}$.

\subsection{Hottel's Crossed-String Algebraic View-Factor Rule \& Cavity Effective Absorptance}
\idxtopic{view factor!Hottel's crossed-string method}
\idxtopic{cavity radiation!hohlraum effective absorptance}
\begin{enumerate}
  \item \textbf{Hottel's 2-D Crossed-String Method (Zero Calculus or Trigonometry Required! Lienhard Table 10.2 Case 4, p.~561 \src{S12}):} For any two infinitely long (2-D) surfaces $A_1$ and $A_2$ (which may be curved, tilted, or partially obstructed), imagine stretching tight strings between their four outer endpoints $(a, b)$ on $A_1$ and $(c, d)$ on $A_2$. The exact view factor $F_{1\to 2}$ equals the sum of the two \textbf{crossed} diagonal string lengths minus the sum of the two \textbf{uncrossed} side string lengths, divided by $2 A_1$:
  \begin{equation}
    \boxed{F_{1\to 2} = \frac{\left(L_{\mathrm{crossed},ad} + L_{\mathrm{crossed},bc}\right) - \left(L_{\mathrm{uncrossed},ac} + L_{\mathrm{uncrossed},bd}\right)}{2 A_1}}
    \label{eq:12:hottel-strings}
  \end{equation}
  For two parallel directly opposed strips of equal width $w$ separated by distance $h$ ($L_{\mathrm{crossed}} = \sqrt{w^2 + h^2},\; L_{\mathrm{uncrossed}} = h$), Eq.~\eref{eq:12:hottel-strings} immediately yields Lienhard Table~10.2 Case~1: $F_{1\to 2} = \sqrt{1 + (h/w)^2} - (h/w)$.
  \item \textbf{Concentric Cylinders and Concentric Spheres ($r_1 < r_2$, Lienhard Tables 10.2--10.3 \src{S12}):} Because the inner convex surface $1$ cannot see itself ($F_{1\to 1} = 0$), all radiation leaving $1$ strikes $2$ ($F_{1\to 2} = 1$). By reciprocity ($A_1 F_{1\to 2} = A_2 F_{2\to 1}$) and summation on the concave outer surface $2$ ($F_{2\to 1} + F_{2\to 2} = 1$):
  \begin{equation}
    F_{1\to 2} = 1, \qquad F_{2\to 1} = \frac{A_1}{A_2} = \begin{cases} r_1 / r_2 & \text{(concentric cylinders)} \\ (r_1 / r_2)^2 & \text{(concentric spheres)} \end{cases}, \qquad F_{2\to 2} = 1 - \frac{A_1}{A_2}
    \label{eq:12:concentric-vf}
  \end{equation}
  \item \textbf{Cavity (\emph{Hohlraum}) Effective Absorptance \& Emittance Enhancement:} When a cavity of interior surface area $A_{\mathrm{cav}}$ and intrinsic wall absorptance $\alpha_s = \varepsilon_s$ opens to the outside through an aperture of area $A_{\mathrm{ap}} \le A_{\mathrm{cav}}$, any ray entering $A_{\mathrm{ap}}$ undergoes an average of $A_{\mathrm{cav}}/A_{\mathrm{ap}}$ internal bounces before escaping, boosting the effective aperture absorptance $\alpha_{\mathrm{cav}}$ to:
  \begin{equation}
    \boxed{\alpha_{\mathrm{cav}} = \varepsilon_{\mathrm{cav}} = \frac{\alpha_s}{\alpha_s + \left(1 - \alpha_s\right)\dfrac{A_{\mathrm{ap}}}{A_{\mathrm{cav}}}} \xrightarrow{A_{\mathrm{ap}} \ll A_{\mathrm{cav}}} 1.000 \quad (\text{ideal blackbody cavity})}
    \label{eq:12:cavity-hohlraum}
  \end{equation}
\end{enumerate}

\subsection{Oppenheim Radiosity-Space Resistance Network and $N$ Radiation Shields}
In a diffuse-gray enclosure, each opaque surface $i$ (area $A_i$, temperature $T_i$, emittance $\varepsilon_i$, reflectance $\rho_i = 1 - \varepsilon_i$) has a surface resistance $R_{\mathrm{surf},i} = \frac{1-\varepsilon_i}{\varepsilon_i A_i}$ connecting its blackbody emissive potential $E_{b,i} = \sigma T_i^4$ to its outgoing \textbf{radiosity node} $J_i = \varepsilon_i\sigma T_i^4 + (1-\varepsilon_i)G_i$ (Lienhard §10.4, pp.~569--573 \src{S12}). Two radiosity nodes $J_1$ and $J_2$ are connected by the geometric space resistance $R_{\mathrm{space},12} = \frac{1}{A_1 F_{1\to 2}} = \frac{1}{A_2 F_{2\to 1}}$, as drawn in \fref{fig:12:shields-viewfactor}.
When $N$ thin parallel reflective sheets (shields $s_1, \dots, s_N$, each with emittance $\varepsilon_s$ on both faces and $F=1$) are placed in series between parallel walls $1$ ($\varepsilon_1$) and $2$ ($\varepsilon_2$):
\begin{equation}
  \dot{Q}_{\mathrm{net},1\to 2} = \frac{\sigma A\left(T_1^4 - T_2^4\right)}{\left(\dfrac{1}{\varepsilon_1} + \dfrac{1}{\varepsilon_2} - 1\right) + \sum_{k=1}^N \left(\dfrac{2}{\varepsilon_{s,k}} - 1\right)} \xrightarrow{\varepsilon_1=\varepsilon_2=\varepsilon_s=\varepsilon} \frac{\sigma A\left(T_1^4 - T_2^4\right)}{(N + 1)\left(\dfrac{2}{\varepsilon} - 1\right)} = \frac{\dot{Q}_{0\text{ shields}}}{N + 1}
  \label{eq:12:n-shields}
\end{equation}
A single floating shield ($N=1$) with $\varepsilon_s = \varepsilon_1 = \varepsilon_2$ reaches steady radiative equilibrium at the fourth-power mean temperature $T_s = \left[\frac{1}{2}(T_1^4 + T_2^4)\right]^{1/4}$ and cuts radiative heat transfer in half; if the shield has much lower emissivity ($\varepsilon_s = 0.04$ aluminum foil) than cardboard walls ($\varepsilon_1 = \varepsilon_2 = 0.92$), a single foil sheet reduces radiative flux by $\frac{2/0.92 - 1 + 2/0.04 - 1}{2/0.92 - 1} = \mathbf{42.7\times}$!

% ---------------- FIGURE 12.2: fig:12:shields-viewfactor (shield network) -----
\begin{center}
\begin{tikzpicture}[font=\footnotesize, >=Latex]
  % Physical plates + shield
  \fill[gray!35] (0.2,0.4) rectangle (0.45,3.1);
  \draw[thick] (0.2,0.4) rectangle (0.45,3.1);
  \node[above, align=center] at (0.32, 3.1) {Wall 1\\($T_1,\varepsilon_1$)};
  \draw[teal, ultra thick] (2.2,0.4) -- (2.2,3.1);
  \node[teal, font=\bfseries\footnotesize, above, align=center] at (2.2, 3.1) {Foil shield\\($T_s,\ \varepsilon_s = 0.04$)};
  \fill[gray!35] (3.95,0.4) rectangle (4.2,3.1);
  \draw[thick] (3.95,0.4) rectangle (4.2,3.1);
  \node[above, align=center] at (4.08, 3.1) {Wall 2\\($T_2,\varepsilon_2$)};
  \draw[->, rust, thick] (0.55,2.4) -- (2.1,2.4);
  \draw[->, teal, thick] (2.1,1.95) -- (0.55,1.95) node[midway, below] {$\rho_s = 0.96$};
  \draw[->, rust, dashed] (2.3,1.2) -- (3.85,1.2) node[midway, above] {$\dot{Q}_{\mathrm{rad}} \ll \dot{Q}_0$};
  % Series resistance chain, two lines
  \node[anchor=west, align=left, draw=navy, fill=boxbg, rounded corners=1.5pt, inner sep=4pt] at (4.9, 2.05)
    {$E_{b1}=\sigma T_1^4 \xrightarrow{\frac{1-\varepsilon_1}{\varepsilon_1 A_1}} J_1 \xrightarrow{\frac{1}{A_1 F_{1s}}} J_{s,1} \xrightarrow{\frac{1-\varepsilon_s}{\varepsilon_s A_s}} E_{bs}=\sigma T_s^4$\\[4pt]
     $\phantom{E_{b1}=\sigma T_1^4} \xrightarrow{\frac{1-\varepsilon_s}{\varepsilon_s A_s}} J_{s,2} \xrightarrow{\frac{1}{A_s F_{s2}}} J_2 \xrightarrow{\frac{1-\varepsilon_2}{\varepsilon_2 A_2}} E_{b2}=\sigma T_2^4$};
  \node[anchor=west, align=left, text=rust, font=\bfseries\footnotesize] at (4.9, 0.35)
    {Equal emittances, $F_{1s} = F_{s2} = 1$: \ $T_s = \left[\tfrac{1}{2}\left(T_1^4 + T_2^4\right)\right]^{1/4}$\\(always closer to the hotter wall)};
\end{tikzpicture}
\captionof{figure}{Oppenheim surface-and-space resistance network with one low-emissivity aluminum-foil radiation shield ($\varepsilon_s = 0.04$) between two walls: six resistances in series \src{S12}.}
\label{fig:12:shields-viewfactor}
\idxfig{radiation shield!Oppenheim resistance network}
\end{center}

% ---------------- FIGURE 12.3: fig:12:lamp-viewfactor (coaxial disks) ---------
\begin{center}
\begin{tikzpicture}[font=\footnotesize, >=Latex]
  % Lamp face at top
  \fill[amber!30, draw=amber, very thick] (1.0, 3.0) rectangle (3.0, 3.25);
  \node[right, font=\bfseries\footnotesize] at (3.1, 3.12) {R40 lamp face, $r_1 = \SI{6.35}{\centi\meter}$, at height $h$};
  % Top aperture of 30 cm cube
  \draw[navy, ultra thick] (0.0, 1.9) -- (4.0, 1.9);
  \node[right, font=\bfseries\footnotesize, text=navy] at (4.1, 1.9) {aperture at $z_{\mathrm{ap}} = \SI{30.0}{\centi\meter}$, $r_2 = \SI{15}{\centi\meter}$};
  % Conical reflector walls tapering to cup mouth
  \draw[teal, ultra thick] (0.0, 1.9) -- (1.4, 0.65);
  \draw[teal, ultra thick] (4.0, 1.9) -- (2.6, 0.65);
  % Cup at bottom
  \fill[navy!15, draw=navy, thick] (1.4, 0.65) -- (1.52, 0.15) -- (2.48, 0.15) -- (2.6, 0.65) -- cycle;
  \node[right, font=\bfseries\footnotesize] at (2.7, 0.4) {9 oz cup, $r_{\mathrm{cup}} = \SI{3.75}{\centi\meter}$};
  % Dimension arrow
  \draw[<->, rust, thick] (-0.25, 1.9) -- (-0.25, 3.0) node[midway, left] {$d_{\mathrm{ap}} = h - 30.0$};
  \node[anchor=north west, fill=white, draw=amber, inner sep=2pt, font=\bfseries\footnotesize, align=left] at (7.2, 1.5)
    {At $h = \SI{40}{\centi\meter}$: $F_{\mathrm{bulb}\to\mathrm{ap}} = 0.665$\\bare cup on the table: $F = 0.0147$\\ratio $= 45.1\times$ (Ex.~\ref{ex:12.2})};
\end{tikzpicture}
\captionof{figure}{Coaxial-disk view-factor geometry between the $\SI{12.7}{\centi\meter}$ R40 lamp face at height $h$ and a $\SI{30}{\centi\meter}$-wide collector aperture at $z_{\mathrm{ap}} = \SI{30.0}{\centi\meter}$, compared with a bare 9~oz cup.}
\label{fig:12:lamp-viewfactor}
\idxfig{view factor!lamp to collector aperture (coaxial disks)}
\end{center}

\subsection{Planetary Radiative Equilibrium and Multi-Layer Atmospheric Greenhouse Model}
\idxtopic{planetary radiative equilibrium!greenhouse model}
A rapidly rotating spherical planet of radius $R_p$ and Bond albedo $A_{\mathrm{alb}}$ at solar constant irradiance $S_0$ $[\si{\watt\per\square\meter}]$ intercepts solar power over its projected disk area $\pi R_p^2$ and radiates thermal IR over its full spherical surface area $4\pi R_p^2$ (tested on the 2026--2027 Cornell Invitational \src{S25} and Lienhard §10.6 pp.~594--601 \src{S12}):
\begin{equation}
  (1 - A_{\mathrm{alb}}) S_0 \left(\pi R_p^2\right) = 4\pi R_p^2 \sigma T_{\mathrm{eq}}^4 \implies T_{\mathrm{eq}} = \left[\frac{(1 - A_{\mathrm{alb}}) S_0}{4\sigma}\right]^{1/4}, \qquad T_{\mathrm{surf}, N\text{ IR-black layers}} = (N + 1)^{1/4}\,T_{\mathrm{eq}}
  \label{eq:12:planetary-eq}
\end{equation}
If a single atmospheric layer ($N=1$) has solar transmittance $\tau_{\mathrm{SW}}$ and thermal-IR absorptance/emittance $\varepsilon_{\mathrm{LW}}$, solving the two-node radiative energy balance gives $T_{\mathrm{surf}} = \left[\frac{(1-A_{\mathrm{alb}})S_0}{4\sigma}\cdot\frac{1 + \tau_{\mathrm{SW}}(1 - \varepsilon_{\mathrm{LW}}/2)}{1 - \varepsilon_{\mathrm{LW}}/2}\right]^{1/4}$, which reduces to $2^{1/4}T_{\mathrm{eq}} = 1.1892\,T_{\mathrm{eq}}$ when $\tau_{\mathrm{SW}}=1$ and $\varepsilon_{\mathrm{LW}}=1$.

% ----------------------------------------------------------------------------
% L4: WORKED EXAMPLES (5 Fully Verified Examples Matching check_sec10_13.py)
% ----------------------------------------------------------------------------
\section{Worked Competition Examples (§12)}
\label{sec:12.3}

\begin{workedexample}[2400 K Heat-Lamp Tungsten Filament: Wien Peak, Exitance \& Lambertian Radiance]{2027 Cornell E1 \src{S25}}{ex:12.1}
The tungsten filament of a $\SI{125.0}{\watt}$ R40 heat lamp operates at $T_{\mathrm{fil}} = \SI{2400}{\kelvin}$ with total hemispherical emittance $\varepsilon_W = 0.350$ (Lienhard Table~10.1 p.~546 \src{S12}). Compute (a) the Wien peak wavelength $\lambda_{\max}$, (b) the hemispherical radiant exitance $M_{\mathrm{fil}}$, (c) the directional Lambertian radiance $L_{\mathrm{fil}}$, and (d) the effective radiating surface area $A_{\mathrm{fil}}$ of the filament assuming negligible base conduction.
\tcblower
\textbf{Solution.} (a) By Wien's displacement law: $\lambda_{\max} = \dfrac{b_{\mathrm{Wien}}}{T_{\mathrm{fil}}} = \dfrac{\SI{2897.772}{\micro\meter\kelvin}}{\SI{2400}{\kelvin}} = \mathbf{\SI{1.2074}{\micro\meter}}$ ($\SI{1207.4}{\nano\meter}$, near-IR).\\
(b) By the Stefan--Boltzmann law:
\[
  M_{\mathrm{fil}} = \varepsilon_W \sigma T_{\mathrm{fil}}^4 = (0.350)(\SI{5.670374e-8}{\watt\per\square\meter\per\kelvin\tothe{4}})(2400)^4 = \mathbf{\SI{6.585e5}{\watt\per\square\meter}}
\]
(c) By Eq.~\eref{eq:12:lambert-radiance}, the Lambertian directional radiance is $L_{\mathrm{fil}} = \dfrac{M_{\mathrm{fil}}}{\pi} = \dfrac{\SI{6.5845e5}{}}{\pi} = \mathbf{\SI{2.096e5}{\watt\per\square\meter\per\steradian}}$.\\
(d) Required filament surface area: $A_{\mathrm{fil}} = \dfrac{\SI{125.0}{\watt}}{\SI{6.5845e5}{\watt\per\square\meter}} = \SI{1.898e-4}{\square\meter} = \mathbf{\SI{1.898}{\square\centi\meter}}$.
\verifiedby{Checked in \texttt{check\_sec10\_13.py} (\texttt{EX-12.1-WIEN}, \texttt{EX-12.1-RADIANCE}): $\lambda_{\max} = \SI{1.2074}{\micro\meter}$, $L_{\mathrm{fil}} = \SI{2.0959e5}{\watt\per\square\meter\per\steradian}$.}
\end{workedexample}

\begin{workedexample}[Coaxial-Disk View Factor: Top Collector Aperture vs.\ Bare 9 oz Cup]{Lienhard Table 10.3 \src{S12}}{ex:12.2}
An R40 heat lamp of face diameter $D_1 = \SI{12.70}{\centi\meter}$ ($r_1 = \SI{6.35}{\centi\meter}$) is mounted at regional minimum height $h = \SI{40.0}{\centi\meter}$ above the Testing Surface (\fref{fig:12:lamp-viewfactor}). Compare the coaxial-disk view factor $F_{1\to 2}$ from the lamp face to: (a) a collector's inscribed circular top aperture of radius $r_2 = \SI{15.0}{\centi\meter}$ elevated to $z_{\mathrm{ap}} = \SI{30.0}{\centi\meter}$ ($d_{\mathrm{ap}} = 40.0 - 30.0 = \SI{10.0}{\centi\meter}$), versus (b) a bare 9~oz cup of mouth radius $r_{\mathrm{cup}} = \SI{3.75}{\centi\meter}$ resting on the Testing Surface with rim height $z_{\mathrm{cup}} = \SI{10.0}{\centi\meter}$ ($d_{\mathrm{cup}} = 40.0 - 10.0 = \SI{30.0}{\centi\meter}$).
\tcblower
\textbf{Solution.} Using Eq.~\eref{eq:12:coaxial-disks}:\\
(a) For the elevated collector top aperture ($d_{\mathrm{ap}} = \SI{10.0}{\centi\meter}$): $R_1 = \frac{6.35}{10.0} = 0.635$, $R_2 = \frac{15.0}{10.0} = 1.500$, so $X = 1 + \frac{1 + 1.500^2}{0.635^2} = 9.0600$, and:
\[
  F_{\mathrm{bulb}\to\mathrm{ap}} = \frac{1}{2}\!\left[9.0600 - \sqrt{9.0600^2 - 4\left(\frac{1.500}{0.635}\right)^{\!2}}\right] = \frac{9.0600 - 7.7307}{2} = \mathbf{0.6647} \quad (\SI{66.5}{\percent}\text{ intercepted!})
\]
(b) For the bare cup on the table ($d_{\mathrm{cup}} = \SI{30.0}{\centi\meter}$): $R_1 = \frac{6.35}{30.0} = 0.21167$, $R_2 = \frac{3.75}{30.0} = 0.12500$, $X = 1 + \frac{1 + 0.125^2}{0.21167^2} = 23.669$, giving $F_{\mathrm{bulb}\to\mathrm{cup}} = \mathbf{0.01474}$ ($\SI{1.47}{\percent}$). Raising and widening the aperture captures $\frac{0.6647}{0.01474} = \mathbf{45.1\times}$ more diffuse lamp flux!
\verifiedby{Checked in \texttt{check\_sec10\_13.py} (\texttt{EX-12.2-F-AP40}, \texttt{EX-12.2-GAIN}): $F_{\mathrm{ap}} = 0.66465$, ratio $= 45.08\times$.}
\end{workedexample}

\begin{workedexample}[Multi-Layer Aluminum Foil Radiation Shields vs.\ Cardboard Walls]{Lienhard Eq.\ 10.31 \src{S12}}{ex:12.3}
The inner cup cradle wall at $T_1 = \SI{45.0}{\celsius}$ ($\SI{318.15}{\kelvin}$) faces the outer collector box wall at $T_2 = \SI{22.0}{\celsius}$ ($\SI{295.15}{\kelvin}$). Compute the radiative heat flux $q''_{\mathrm{rad}}$ between the walls for: (a) bare corrugated cardboard walls ($\varepsilon_1 = \varepsilon_2 = 0.92$), (b) lining both walls with household aluminum foil ($\varepsilon_1 = \varepsilon_2 = 0.04,\; N = 0$ intermediate shields), and (c) inserting $N = 2$ floating double-sided aluminum-foil radiation shields ($\varepsilon_s = 0.04$) between the foil-lined walls. Also find the equilibrium temperature $T_s$ of a single floating shield ($N=1$).
\tcblower
\textbf{Solution.} Blackbody potential difference is $\sigma(T_1^4 - T_2^4) = (\SI{5.670374e-8}{})(318.15^4 - 295.15^4) = \SI{150.64}{\watt\per\square\meter}$. By Eq.~\eref{eq:12:n-shields}:
{\footnotesize\begin{align*}
  \text{(a) Bare cardboard ($\varepsilon=0.92,\; N=0$):}\ \mathcal{F}_{12} &= \frac{1}{2/0.92 - 1} = 0.85185 \implies q''_{\mathrm{card}} = (0.85185)(150.64) = \mathbf{\SI{128.3}{\watt\per\square\meter}} \\
  \text{(b) Foil-lined walls ($\varepsilon=0.04,\; N=0$):}\ \mathcal{F}_{12} &= \frac{1}{2/0.04 - 1} = \frac{1}{49} = 0.02041 \implies q''_{\mathrm{foil},0} = \frac{150.64}{49} = \mathbf{\SI{3.074}{\watt\per\square\meter}}\ (41.7\times\text{ drop}) \\
  \text{(c) Foil walls + $N=2$ shields:}\ \mathcal{F}_{12} &= \frac{1}{(2+1)(49)} = \frac{1}{147} \implies q''_{\mathrm{foil},2} = \frac{150.64}{147} = \mathbf{\SI{1.025}{\watt\per\square\meter}}\ (\mathbf{125.2\times}\text{ drop!})
\end{align*}}
For $N=1$ equal-$\varepsilon$ shield, $T_s = \left[\frac{1}{2}(318.15^4 + 295.15^4)\right]^{1/4} = \SI{307.30}{\kelvin} = \mathbf{\SI{34.15}{\celsius}}$.
\verifiedby{Checked in \texttt{check\_sec10\_13.py} (\texttt{EX-12.3-Q-CARD}, \texttt{EX-12.3-Q-SHIELD}): $q''_{\mathrm{card}} = \SI{128.32}{\watt\per\square\meter}$, $q''_{\mathrm{foil},2} = \SI{1.025}{\watt\per\square\meter}$.}
\end{workedexample}

\begin{workedexample}[Two-Band Selective Absorber vs.\ Non-Selective Black and White Paint]{Lienhard Ex.\ 10.12 \src{S12}}{ex:12.4}
An insulated flat collector plate receives short-wave irradiance $G = \SI{600.0}{\watt\per\square\meter}$ in ambient air at $T_{\mathrm{air}} = \SI{35.0}{\celsius}$ ($\SI{308.15}{\kelvin}$) with $h_{\mathrm{conv}} = \SI{8.00}{\watt\per\square\meter\per\kelvin}$ and effective radiative surround temperature $T_{\mathrm{sky}} = \SI{18.0}{\celsius}$ ($\SI{291.15}{\kelvin}$). Using Eq.~\eref{eq:12:selective-balance} and \tref{tab:12:selective-emissivity}, find the steady equilibrium plate temperature $T_{\mathrm{plate}}$ for: (a) non-selective black paint ($\alpha_{\mathrm{SW}} = 0.90,\; \varepsilon_{\mathrm{MIR}} = 0.90$), (b) white acrylic paint ($\alpha_{\mathrm{SW}} = 0.26,\; \varepsilon_{\mathrm{MIR}} = 0.90$), and (c) a selective Black-Chrome-on-Nickel absorber ($\alpha_{\mathrm{SW}} = 0.95,\; \varepsilon_{\mathrm{MIR}} = 0.09$).
\tcblower
\textbf{Solution.} Setting $q''_{\mathrm{net}} = 0$ in Eq.~\eref{eq:12:selective-balance}: $8.00(T_p - 308.15) + \varepsilon_{\mathrm{MIR}}(\SI{5.6704e-8}{})(T_p^4 - 291.15^4) = 600.0\,\alpha_{\mathrm{SW}}$.\\
(a) Non-selective black ($\alpha=\varepsilon=0.90$, absorbed flux $\SI{540}{\watt\per\square\meter}$): $T_p = \SI{338.12}{\kelvin} = \mathbf{\SI{64.97}{\celsius}}$ ($\approx\SI{65}{\celsius}$ in Lienhard Ex.~10.12).\\
(b) White acrylic paint ($\alpha=0.26,\; \varepsilon=0.90$, absorbed flux $\SI{156}{\watt\per\square\meter}$): $T_p = \SI{312.59}{\kelvin} = \mathbf{\SI{39.44}{\celsius}}$ ($\SI{25.5}{\celsius}$ cooler than black!).\\
(c) Selective Black Chrome ($\alpha=0.95,\; \varepsilon=0.09$, absorbed flux $\SI{570}{\watt\per\square\meter}$ with $10\times$ lower mid-IR re-radiation!):
\[
  8.00(T_p - 308.15) + (0.09)(\SI{5.6704e-8}{})(T_p^4 - 291.15^4) = 570.0 \implies T_p = \SI{371.79}{\kelvin} = \mathbf{\SI{98.64}{\celsius}}
\]
Suppressing $\varepsilon_{\mathrm{MIR}}$ from $0.90$ to $0.09$ boosts the stagnation temperature by $+\SI{33.7}{\celsius}$!
\verifiedby{Brent root-finding in \texttt{check\_sec10\_13.py} (\texttt{EX-12.4-BLACK}, \texttt{EX-12.4-SELECTIVE}): $T_{\mathrm{black}} = \SI{64.966}{\celsius}$, $T_{\mathrm{sel}} = \SI{98.645}{\celsius}$.}
\end{workedexample}

\begin{workedexample}[Planetary Radiative Equilibrium \& Single-Layer Greenhouse Sheet]{2027 Cornell E1 \src{S25}}{ex:12.5}
Earth receives top-of-atmosphere solar irradiance $S_0 = \SI{1361}{\watt\per\square\meter}$ (Lienhard §10.6 p.~594 \src{S12}) with global Bond albedo $A_{\mathrm{alb}} = 0.300$. Compute (a) Earth's bare blackbody radiative equilibrium temperature $T_{\mathrm{eq}}$ without an atmosphere, and (b) the surface temperature $T_s$ under an idealized 1-layer greenhouse atmosphere that is transparent to solar radiation ($\tau_{\mathrm{SW}}=1$) and black to thermal infrared ($\varepsilon_{\mathrm{LW}}=1$).
\tcblower
\textbf{Solution.} (a) By Eq.~\eref{eq:12:planetary-eq}, the globally averaged absorbed solar flux is $q''_{\mathrm{abs}} = \frac{1}{4}(1 - 0.300)(1361) = \SI{238.175}{\watt\per\square\meter}$:
\[
  T_{\mathrm{eq}} = \left(\frac{238.175}{\SI{5.670374e-8}{}}\right)^{\!1/4} = \mathbf{\SI{254.58}{\kelvin}} = -\SI{18.57}{\celsius}
\]
(b) With $N=1$ IR-absorbing atmospheric sheet radiating $\sigma T_{\mathrm{atm}}^4 = \sigma T_{\mathrm{eq}}^4$ both upward to space and downward to the ground, the ground absorbs $2\sigma T_{\mathrm{eq}}^4$, giving $T_s = 2^{1/4}T_{\mathrm{eq}} = (1.18921)(254.58) = \mathbf{\SI{302.75}{\kelvin}} = +\SI{29.60}{\celsius}$.
\verifiedby{Checked in \texttt{check\_sec10\_13.py} (\texttt{EX-12.5-TEQ}, \texttt{EX-12.5-TGH}): $T_{\mathrm{eq}} = \SI{254.578}{\kelvin}$, $T_s = \SI{302.746}{\kelvin}$.}
\end{workedexample}

\begin{workedexample}[Hottel Crossed-String View Factor \& Cavity Effective Absorptance]{Lienhard Table 10.2 \src{S12}}{ex:12.6}
(a) Use Hottel's crossed-string rule Eq.~\eref{eq:12:hottel-strings} to find the view factor $F_{1\to 2}$ between two parallel, directly opposed long strips of width $w = \SI{10.0}{\centi\meter}$ separated by distance $h = \SI{10.0}{\centi\meter}$. (b) A 2027 cup-cradle cavity has top aperture area $A_{\mathrm{ap}} = \SI{44.18}{\square\centi\meter}$ (matching the 9~oz cup mouth) and interior absorbing surface area $A_{\mathrm{cav}} = \SI{220.9}{\square\centi\meter}$ ($A_{\mathrm{ap}}/A_{\mathrm{cav}} = 0.200$). If the interior surface has intrinsic single-bounce absorptance $\alpha_s = 0.750$, compute the effective cavity absorptance $\alpha_{\mathrm{cav}}$.
\tcblower
\textbf{Solution.} (a) Uncrossed side strings have length $L_{u1} = L_{u2} = h = \SI{0.100}{\meter}$, while crossed diagonal strings have Pythagorean length $L_{c1} = L_{c2} = \sqrt{w^2 + h^2} = \sqrt{0.100^2 + 0.100^2} = \SI{0.141421}{\meter}$. By Eq.~\eref{eq:12:hottel-strings}:
\[
  F_{1\to 2} = \frac{2(0.141421) - 2(0.100)}{2(0.100)} = \sqrt{2} - 1 = \mathbf{0.4142}
\]
(b) By Eq.~\eref{eq:12:cavity-hohlraum}, cavity multiple-reflection trapping elevates $\alpha_{\mathrm{cav}}$ from $0.750$ to:
\[
  \alpha_{\mathrm{cav}} = \frac{0.750}{0.750 + (1 - 0.750)(0.200)} = \frac{0.750}{0.800} = \mathbf{0.9375} \quad (+\SI{18.75}{\percent}\text{ cavity trapping gain!})
\]
\verifiedby{Checked in \texttt{check\_sec10\_13.py} (\texttt{EX-12.6-CAV}, \texttt{EX-12.6-HOTTEL}): $\alpha_{\mathrm{cav}} = 0.9375$, $F_{1\to 2} = 0.414214$.}
\end{workedexample}

\begin{workedexample}[Linearized Radiation Coefficient $h_{\mathrm{rad}}$ vs.\ Exact Stefan--Boltzmann Power]{Lienhard Ex.\ 2.7 \src{S12}}{ex:12.7}
A small electronic component of surface area $A = \SI{1.3345e-4}{\square\meter}$ and emittance $\varepsilon = 0.900$ at $T_1 = \SI{72.30}{\celsius}$ ($\SI{345.45}{\kelvin}$) radiates inside a large cabinet at $T_2 = \SI{35.00}{\celsius}$ ($\SI{308.15}{\kelvin}$). Compute (a) the mean temperature $T_m$ and linearized radiation coefficient $h_{\mathrm{rad}} = 4\varepsilon\sigma T_m^3$, (b) the linearized radiative heat loss $\dot{Q}_{\mathrm{rad,lin}} = h_{\mathrm{rad}}A(T_1 - T_2)$ versus exact $\dot{Q}_{\mathrm{rad,ex}} = \varepsilon\sigma A(T_1^4 - T_2^4)$, and (c) verify that the fractional discrepancy matches the Taylor remainder $\frac{1}{4}(\Delta T/T_m)^2$ (Lienhard Eq.~2.28, p.~72 \src{S12}).
\tcblower
\textbf{Solution.} (a) $T_m = \frac{1}{2}(345.45 + 308.15) = \SI{326.80}{\kelvin}$, and $\Delta T = \SI{37.30}{\kelvin}$:
\[
  h_{\mathrm{rad}} = 4(0.900)(\SI{5.670374e-8}{})(326.80)^3 = \mathbf{\SI{7.125}{\watt\per\square\meter\per\kelvin}}
\]
(b) Linearized power: $\dot{Q}_{\mathrm{rad,lin}} = (7.1246)(\SI{1.3345e-4}{})(37.30) = \mathbf{\SI{0.03546}{\watt}}$, compared to exact $\dot{Q}_{\mathrm{rad,ex}} = (0.900)(\SI{5.670374e-8}{})(\SI{1.3345e-4}{})(345.45^4 - 308.15^4) = \mathbf{\SI{0.03558}{\watt}}$.\\
(c) Fractional difference is $\frac{0.035580 - 0.035464}{0.035580} = \mathbf{0.00325} = \SI{0.325}{\percent}$, matching $\frac{1}{4}\!\left(\frac{37.30}{326.80}\right)^{\!2} = \mathbf{0.003257}$!
\verifiedby{Checked in \texttt{check\_sec10\_13.py} (\texttt{EX-12.7-HRAD}, \texttt{EX-12.7-EXACT}): $h_{\mathrm{rad}} = \SI{7.1246}{\watt\per\square\meter\per\kelvin}$.}
\end{workedexample}

% ----------------------------------------------------------------------------
% L5: TRAPS & DIAGNOSTICS + L6: 2027 DEVICE APPLICATION
% ----------------------------------------------------------------------------
\section{Radiation Traps, Diagnostics \& 2027 Optical Concentrator Rules}
\label{sec:12.4}

\begin{trapbox}[Top Thermal Radiation Traps on Science Olympiad Exams]
\idxerr{radiation!factoring T1^4 - T2^4 as (T1 - T2)^4}
\idxerr{radiation!using 2*pi instead of pi in Lambertian radiance}
\idxerr{radiation!assuming white paint has low mid-IR emissivity}
\par\smallskip\noindent
\begin{errtab}
\textbf{1. Writing $(T_1 - T_2)^4$ instead of $(T_1^4 - T_2^4)$, or leaving $T$ in Celsius}
  & $318.15^4 - 295.15^4 = \num{2.659e9}\ \si{\kelvin\tothe{4}}$, whereas $(45 - 22)^4 = \num{2.798e5}$ (off by $\mathbf{9500\times}$!).
  & Convert $T_1, T_2$ to Kelvin first, raise each to the 4th power, then subtract. \\
\textbf{2. Dividing hemispherical exitance $M = \sigma T^4$ by $2\pi\ \si{\steradian}$ instead of $\pi\ \si{\steradian}$ to get radiance $L$}
  & Off by a factor of $2\times$; forgets the Lambertian $\cos\theta$ projected-area factor in $\int_0^{\pi/2}\cos\theta\sin\theta\,d\theta = 1/2$.
  & Diffuse radiance is always $L = M/\pi = \varepsilon\sigma T^4/\pi$ \src{S12,S25}. \\
\textbf{3. Assuming visually white surfaces (white paper/paint) have low mid-IR emissivity $\varepsilon_{\mathrm{MIR}}$}
  & White paper and white acrylic paint have $\varepsilon_{\mathrm{MIR}} = \mathbf{0.90\text{--}0.96}$ at $\lambda \approx \SI{9.6}{\micro\meter}$ — they radiate thermal IR like a blackbody!
  & Only polished metals ($\mathrm{Al}, \mathrm{Ag}, \mathrm{Cu}, \mathrm{Au}$) have low $\varepsilon_{\mathrm{MIR}} \approx 0.02\text{--}0.05$. \\
\textbf{4. Assuming view factor reciprocity is $F_{1\to 2} = F_{2\to 1}$ when $A_1 \ne A_2$}
  & Violates Second Law! Reciprocity includes areas: $A_1 F_{1\to 2} = A_2 F_{2\to 1}$, so $F_{1\to 2} = F_{2\to 1}$ \emph{only} when $A_1 = A_2$.
  & Small body $1$ enclosed by large body $2$ has $F_{1\to 2} = 1$, while $F_{2\to 1} = A_1/A_2 \ll 1$. \\
\end{errtab}
\end{trapbox}

\begin{devicebox}[2027 Heat-Lamp Collector Optical \& Radiative Design Rules \src{S01,S02,S12,S19}]
\idxdev{aluminum foil!shiny vs dull specular reflectance}
\idxdev{greenhouse glazing!BoPET vs glass cutoff at 2.7 um}
\begin{enumerate}
  \item \textbf{Shiny vs.\ Dull Aluminum Foil — Specular Concentrator vs.\ Radiant Shield:} Both faces of household aluminum foil have virtually identical total hemispherical IR reflectance ($\rho_{\mathrm{NIR,tot}} = 0.95\text{--}0.97$) and ultra-low mid-IR emittance ($\varepsilon_{\mathrm{MIR}} = 0.030\text{--}0.045$), performing identically inside insulation air gaps. However, inside the upper concentrating funnel, \textbf{always orient the SHINY side toward the heat lamp} because its specular (mirror-like) reflectance is $\rho_{\mathrm{spec}} \approx 0.88\text{--}0.91$ (versus $\sim 0.62$ on the micro-matte dull side), focusing directional rays straight into the 9~oz cup rather than scattering rays back out the top aperture!
  \item \textbf{Two-Band Greenhouse Trapping with Removable Clear BoPET Film (FAQ 10/21/2026 on items on the cup):} The $\SI{2400}{\kelvin}$ tungsten heat lamp emits $>80\%$ of its glass-transmitted power at $\lambda \in [\SI{0.4}{\micro\meter}, \SI{2.7}{\micro\meter}]$, whereas the $\SI{310}{\kelvin}$ water emits at $\lambda \in [\SI{5}{\micro\meter}, \SI{25}{\micro\meter}]$. A $\SI{25}{\micro\meter}$ clear BoPET film resting on the cup rim transmits $\tau_{\mathrm{NIR}} \approx 0.90$ of the incoming lamp beam while blocking mid-IR re-radiation ($\tau_{\mathrm{MIR}} \approx 0.20$) AND halting evaporative heat loss!
\end{enumerate}
\end{devicebox}

% ----------------------------------------------------------------------------
% RAPID COMPETITION PROBLEM & SOLUTION BANK (§12)
% ----------------------------------------------------------------------------
\section{Rapid Competition Problem \& Solution Bank (§12)}
\label{sec:12.5}

\begin{problem}{P12.1}{Concentric Spherical Dewar View Factors \& Self-View Factor $F_{2\to 2}$}{12.2}{5.c.iii, 5.c.vi}{State / National}{60 s}{Reciprocity + Summation}{setting $F_{2\to 2} = 0$ on a concave outer sphere}
A spherical inner vessel of radius $r_1 = \SI{6.00}{\centi\meter}$ ($\varepsilon_1 = 0.050$) is enclosed concentrically inside an outer spherical shell of inner radius $r_2 = \SI{12.0}{\centi\meter}$ ($\varepsilon_2 = 0.050$) across an evacuated vacuum gap. Compute (a) all four view factors $F_{1\to 1}, F_{1\to 2}, F_{2\to 1}, F_{2\to 2}$, and (b) the gray transfer factor $\mathcal{F}_{1\to 2}$.
\end{problem}
\begin{solution}{P12.1}
(a) Because sphere $1$ is convex, $F_{1\to 1} = \mathbf{0}$ and $F_{1\to 2} = \mathbf{1.000}$. By reciprocity Eq.~\eref{eq:12:concentric-vf}, $F_{2\to 1} = \dfrac{A_1}{A_2} = \left(\dfrac{r_1}{r_2}\right)^{\!2} = \left(\dfrac{6.00}{12.0}\right)^{\!2} = \mathbf{0.2500}$. By summation on the concave outer sphere $2$, $F_{2\to 2} = 1 - F_{2\to 1} = \mathbf{0.7500}$ ($75\%$ of diffuse rays leaving the outer sphere strike the outer sphere directly!).\\
(b) By Lienhard Eq.~(10.27): $\mathcal{F}_{1\to 2} = \left[\dfrac{1}{\varepsilon_1} + \dfrac{A_1}{A_2}\!\left(\dfrac{1}{\varepsilon_2} - 1\right)\right]^{-1} = \left[\dfrac{1}{0.050} + 0.250(19.0)\right]^{-1} = \dfrac{1}{24.75} = \mathbf{0.04040}$.
\end{solution}

\begin{problem}{P12.2}{Wien Peak Shift \& Stefan--Boltzmann Scaling Between 100 W and 125 W Lamps}{12.1}{5.c.vi}{State / National}{60 s}{Fourth-Power Scaling}{confusing $\lambda_{\max} \propto T^{-1}$ with $P_{\mathrm{rad}} \propto T^4$}
When an incandescent tungsten filament of constant surface area $A_{\mathrm{fil}}$ and constant emittance $\varepsilon_W$ has its electrical power reduced from $P_1 = \SI{125.0}{\watt}$ (at filament temperature $T_1 = \SI{2400}{\kelvin}$) to $P_2 = \SI{100.0}{\watt}$ (neglecting stem conduction), find the new filament temperature $T_2$ and the new Wien peak wavelength $\lambda_{\max,2}$.
\end{problem}
\begin{solution}{P12.2}
Since $P_{\mathrm{rad}} = \varepsilon_W \sigma A_{\mathrm{fil}} T^4 \propto T^4$, the new filament temperature is:
\[
  T_2 = T_1\left(\frac{P_2}{P_1}\right)^{\!1/4} = (\SI{2400}{\kelvin})\left(\frac{100.0}{125.0}\right)^{\!1/4} = 2400(0.8000)^{1/4} = 2400(0.94574) = \mathbf{\SI{2270}{\kelvin}}
\]
The Wien peak shifts to a slightly longer near-IR wavelength: $\lambda_{\max,2} = \dfrac{\SI{2897.77}{\micro\meter\kelvin}}{\SI{2269.8}{\kelvin}} = \mathbf{\SI{1.277}{\micro\meter}}$ (up from $\SI{1.207}{\micro\meter}$, increasing direct top-millimeter water absorption by $+2.3\%$!).
\end{solution}

\begin{problem}{P12.3}{Asymmetric Radiation Shield Between High-Emissivity and Low-Emissivity Walls}{12.2}{5.c.iii, 5.c.vi}{National / MIT}{75 s}{Oppenheim Network}{assuming shield temperature $T_s^4 = (T_1^4+T_2^4)/2$ when $\varepsilon_1 \ne \varepsilon_2$}
Two large parallel plates at $T_1 = \SI{400.0}{\kelvin}$ ($\varepsilon_1 = 0.800$) and $T_2 = \SI{300.0}{\kelvin}$ ($\varepsilon_2 = 0.800$) are separated by a thin floating sheet that is painted black on Side~A facing Plate~1 ($\varepsilon_{s,A} = 0.800$) and polished aluminum on Side~B facing Plate~2 ($\varepsilon_{s,B} = 0.050$). Find (a) the steady temperature $T_s$ of the asymmetric shield, and (b) does flipping the shield ($180^\circ$) change the net heat flux $q''_{1\to 2}$ or the shield temperature $T_s$?
\end{problem}
\begin{solution}{P12.3}
(a) The resistance between Plate~1 and Shield Side~A is $R_{1s}'' = \frac{1}{0.80} + \frac{1}{0.80} - 1 = \mathbf{1.500}$, while the resistance between Shield Side~B and Plate~2 is $R_{s2}'' = \frac{1}{0.050} + \frac{1}{0.80} - 1 = \mathbf{20.250}$ ($13.5\times$ larger!). At steady state, $q''_{1\to s} = q''_{s\to 2} \implies \frac{\sigma(T_1^4 - T_s^4)}{1.500} = \frac{\sigma(T_s^4 - T_2^4)}{20.250}$:
\[
  T_s^4 = \frac{20.250\,T_1^4 + 1.500\,T_2^4}{20.250 + 1.500} = \frac{20.250(400^4) + 1.500(300^4)}{21.750} = \num{2.4393e10}\ \si{\kelvin\tothe{4}} \implies T_s = \mathbf{\SI{395.2}{\kelvin}}
\]
(b) Flipping the shield leaves the \textbf{total series resistance} $R_{\mathrm{tot}}'' = 1.500 + 20.250 = 21.750$ and \textbf{net heat flux $q''_{1\to 2}$ completely unchanged}, but shifts the shield temperature $T_s$ from $\SI{395.2}{\kelvin}$ (close to $T_1$) to $\SI{310.6}{\kelvin}$ (close to $T_2$)!
\end{solution}

\begin{problem}{P12.4}{Thermodynamic Second-Law Limit on Optical Concentration Ratio $C_{\max}$}{12.2}{5.c.iv, 5.c.vi}{National}{60 s}{Second Law + Reciprocity}{exceeding $1/\sin^2\theta_s$ or $T_{\mathrm{target}} > T_{\mathrm{source}}$}
An R40 heat lamp of radiating disk radius $r_{\mathrm{source}} = \SI{6.35}{\centi\meter}$ at temperature $T_{\mathrm{source}} = \SI{2400}{\kelvin}$ sits at distance $d = \SI{20.0}{\centi\meter}$ above a collector aperture ($\sin\theta_s = r_{\mathrm{source}}/\sqrt{r_{\mathrm{source}}^2 + d^2} = 0.3026$). According to the Second Law of Thermodynamics (conservation of optical \'etendue / view-factor reciprocity $A_1 F_{1\to 2} = A_2 F_{2\to 1}$ with $F_{2\to 1} \le 1$), what is the maximum theoretical geometric concentration ratio $C_{\max} = A_{\mathrm{ap}}/A_{\mathrm{receiver}}$ achievable by any passive 3-D ideal concentrator?
\end{problem}
\begin{solution}{P12.4}
For a 3-D axisymmetric source subtending half-angle $\theta_s$ at the aperture, view-factor reciprocity requires $A_{\mathrm{ap}} F_{\mathrm{ap}\to\mathrm{source}} = A_{\mathrm{ap}}\sin^2\theta_s \le A_{\mathrm{rec}} F_{\mathrm{rec}\to\mathrm{ap}} \le A_{\mathrm{rec}}(1)$. If $C_{\mathrm{geom}} > 1/\sin^2\theta_s$, an isolated receiver could reach a steady temperature $T_{\mathrm{rec}} > T_{\mathrm{source}}$, violating the Clausius statement of the Second Law! Thus:
\[
  C_{\max,\mathrm{3D}} = \frac{1}{\sin^2\theta_s} = \frac{r_{\mathrm{source}}^2 + d^2}{r_{\mathrm{source}}^2} = 1 + \left(\frac{20.0}{6.35}\right)^{\!2} = \mathbf{10.92\times} \quad (\text{and at } d = \SI{35.0}{\centi\meter},\ C_{\max,\mathrm{3D}} = \mathbf{31.38\times})
\]
\end{solution}

\par\smallskip\noindent
\begin{minipage}{\linewidth}
\captionof{table}{Universal Planck Blackbody Fractional Function $F_{0\to\lambda T} \equiv \frac{1}{\sigma T^4}\int_0^{\lambda} M_{\lambda,b}(\lambda',T)\,d\lambda'$ as a Function of $\lambda T$ $[\si{\micro\meter\kelvin}]$ (Lienhard Table~10.2, p.~556 \src{S12,S25}). Fraction of total blackbody power emitted in band $[\lambda_1, \lambda_2]$ is $\Delta F = F_{0\to\lambda_2 T} - F_{0\to\lambda_1 T}$.}
\label{tab:12:blackbody-fraction}
\idxtopic{Planck blackbody fractional function!universal lookup table}
\idxtopic{spectral band fraction!F(0 to lambda*T)}
\footnotesize
\setlength{\tabcolsep}{4pt}
\renewcommand{\arraystretch}{1.12}
\begin{tabularx}{\linewidth}{@{}r r r r r r r X@{}}
\toprule
$\lambda T$ $[\si{\micro\meter\kelvin}]$ & $F_{0\to\lambda T}$ & $\lambda T$ $[\si{\micro\meter\kelvin}]$ & $F_{0\to\lambda T}$ & $\lambda T$ $[\si{\micro\meter\kelvin}]$ & $F_{0\to\lambda T}$ & $\lambda T$ $[\si{\micro\meter\kelvin}]$ & $F_{0\to\lambda T}$ \\
\midrule
$1000$ & $0.00032$ & $2400$ & $0.14026$ & $3800$ & $0.44338$ & $6480$ & $0.77492$ \\
$1400$ & $0.00779$ & $2600$ & $0.18312$ & $4000$ & $0.48088$ & $7000$ & $0.80811$ \\
$1680$ & $0.02661$ & $2880$ & $0.24603$ & $4107.2$ & $0.50000$ & $8000$ & $0.85629$ \\
$1800$ & $0.03934$ & $2897.8$ & $0.25006$ & $4600$ & $0.57928$ & $10000$ & $0.91420$ \\
$2000$ & $0.06673$ & $3200$ & $0.31810$ & $5000$ & $0.63375$ & $15000$ & $0.96898$ \\
$2200$ & $0.10089$ & $3600$ & $0.40361$ & $6000$ & $0.73782$ & $25000$ & $0.99222$ \\
\bottomrule
\end{tabularx}
\end{minipage}
\par\smallskip

\begin{problem}{P12.5}{Planck Band-Fraction Split of a $\SI{2400}{\kelvin}$ Heat Lamp Across Visible, Water-Penetrating Near-IR, and Glass Cutoff}{12.1}{5.c.vi}{National / MIT}{75 s}{Table 12.2 Band Lookup}{integrating Planck law by hand instead of subtracting $F_{0\to\lambda T}$}
Using \tref{tab:12:blackbody-fraction}, determine for a $\SI{2400}{\kelvin}$ blackbody filament: (a) the fraction of power emitted at $\lambda \le \SI{0.700}{\micro\meter}$ ($\lambda T = \SI{1680}{\micro\meter\kelvin}$, visible + UV), (b) the fraction emitted in the water-surface-absorbed near-IR window $\lambda \in [\SI{1.200}{\micro\meter}, \SI{2.700}{\micro\meter}]$ ($\lambda T \in [\SI{2880}{}, \SI{6480}{\micro\meter\kelvin}]$), and (c) the median wavelength $\lambda_{\mathrm{med}}$ that divides the emitted power into two equal $50\%$ halves.
\end{problem}
\begin{solution}{P12.5}
\begin{enumerate}
  \item[(a)] At $\lambda_1 = \SI{0.700}{\micro\meter}$, $\lambda_1 T = (0.700)(2400) = \SI{1680}{\micro\meter\kelvin} \implies F_{0\to 1680} = \mathbf{0.02661}$ ($\mathbf{\SI{2.66}{\percent}}$ visible/UV).
  \item[(b)] At $\lambda_2 = \SI{1.200}{\micro\meter}$, $\lambda_2 T = \SI{2880}{\micro\meter\kelvin} \implies F_{0\to 2880} = 0.24603$. At $\lambda_3 = \SI{2.700}{\micro\meter}$ (bulb glass cutoff), $\lambda_3 T = \SI{6480}{\micro\meter\kelvin} \implies F_{0\to 6480} = 0.77492$. The power fraction in $[\SI{1.20}{\micro\meter}, \SI{2.70}{\micro\meter}]$ is:
  \[
    \Delta F_{1.2\to 2.7\,\mu\mathrm{m}} = F_{0\to 6480} - F_{0\to 2880} = 0.77492 - 0.24603 = \mathbf{0.52889} \quad (\mathbf{\SI{52.89}{\percent}}\text{ of total filament emission!})
  \]
  \item[(c)] From \tref{tab:12:blackbody-fraction}, $F_{0\to\lambda_{\mathrm{med}} T} = 0.50000$ occurs at $\lambda_{\mathrm{med}} T = \SI{4107.2}{\micro\meter\kelvin} \implies \lambda_{\mathrm{med}} = \dfrac{4107.2}{2400} = \mathbf{\SI{1.7113}{\micro\meter}}$ (notice that $\lambda_{\mathrm{med}} = 1.4174\,\lambda_{\max}$ because the Planck tail is skewed toward long wavelengths, while $25.0\%$ of power lies below $\lambda_{\max}$ and $75.0\%$ lies above $\lambda_{\max}$!).
\end{enumerate}
\end{solution}

\begin{problem}{P12.6}{Combined Multimode Convection + Radiation From an Uninsulated vs.\ Foil-Wrapped Pipe}{12.3}{5.c.iii, 5.c.vi}{State / National}{75 s}{Parallel Conductances $h_{\mathrm{tot}} = h_{\mathrm{conv}} + h_{\mathrm{rad}}$}{forgetting radiation equals or exceeds natural convection in still air}
A horizontal oxidized steel steam pipe of outer diameter $D = \SI{50.0}{\milli\meter}$, length $\ell = \SI{2.00}{\meter}$ ($A_s = \pi D\ell = \SI{0.31416}{\square\meter}$), and surface temperature $T_s = \SI{107.0}{\celsius}$ ($\SI{380.15}{\kelvin}$) passes through a large room with still air and walls at $T_\infty = T_{\mathrm{surr}} = \SI{27.0}{\celsius}$ ($\SI{300.15}{\kelvin}$), where the natural-convection heat transfer coefficient is $h_{\mathrm{conv}} = \SI{6.80}{\watt\per\square\meter\per\kelvin}$. Calculate (a) the linearized radiative coefficient $h_{\mathrm{rad}}$ and total heat loss $\dot{Q}_{\mathrm{tot}}$ when the bare oxidized steel has $\varepsilon_1 = 0.850$, and (b) the percentage reduction in total heat loss if a thin layer of shiny aluminum foil tape ($\varepsilon_2 = 0.040$) is wrapped tightly over the pipe without adding any foam thickness.
\end{problem}
\begin{solution}{P12.6}
\begin{enumerate}
  \item[(a)] Using the linearized radiation coefficient at $T_s = \SI{380.15}{\kelvin}$ and $T_{\mathrm{surr}} = \SI{300.15}{\kelvin}$ ($\Delta T = \SI{80.0}{\kelvin}$):
  \[
    h_{\mathrm{rad},1} = \varepsilon_1 \sigma\left(T_s^2 + T_{\mathrm{surr}}^2\right)\left(T_s + T_{\mathrm{surr}}\right) = (0.850)(\num{5.670374e-8})\left(380.15^2 + 300.15^2\right)(680.30) = \mathbf{\SI{7.692}{\watt\per\square\meter\per\kelvin}}
  \]
  Notice that in still air, $h_{\mathrm{rad},1} = \SI{7.69}{\watt\per\square\meter\per\kelvin} > h_{\mathrm{conv}} = \SI{6.80}{\watt\per\square\meter\per\kelvin}$! Total parallel loss is:
  \[
    \dot{Q}_{\mathrm{tot},1} = \left(h_{\mathrm{conv}} + h_{\mathrm{rad},1}\right)A_s\Delta T = (6.80 + 7.692)(0.31416)(80.0) = \mathbf{\SI{364.2}{\watt}}
  \]
  \item[(b)] Wrapping shiny aluminum foil tape reduces $\varepsilon$ from $0.850$ to $0.040$:
  \[
    h_{\mathrm{rad},2} = \frac{0.040}{0.850}(\SI{7.692}{\watt\per\square\meter\per\kelvin}) = \mathbf{\SI{0.362}{\watt\per\square\meter\per\kelvin}}, \qquad
    \dot{Q}_{\mathrm{tot},2} = (6.80 + 0.362)(0.31416)(80.0) = \mathbf{\SI{180.0}{\watt}}
  \]
  This achieves a $\mathbf{\SI{50.6}{\percent}}$ reduction in total heat loss with zero foam thickness!
\end{enumerate}
\end{solution}

\begin{problem}{P12.7}{Thermocouple Bead Radiative Error inside a Cool-Wall Hot-Gas Duct}{12.3}{5.c.iii, 5.c.vi}{National / MIT}{75 s}{Convection = Radiation Balance}{forgetting that thermocouple reads its own bead temperature $T_{tc}$, not gas $T_g$}
An unshielded spherical thermocouple junction ($\varepsilon_{tc} = 0.850$) is inserted into a hot air stream flowing through a large sheet-metal duct whose walls are at $T_w = \SI{350.0}{\kelvin}$ ($\SI{76.85}{\celsius}$). The convective heat transfer coefficient on the small thermocouple bead is $h_{\mathrm{conv}} = \SI{65.0}{\watt\per\square\meter\per\kelvin}$, and the thermocouple reads a steady temperature $T_{tc} = \SI{450.0}{\kelvin}$ ($\SI{176.85}{\celsius}$). Calculate (a) the true hot-gas temperature $T_g$, and (b) the new thermocouple reading if a polished cylindrical silver radiation shield ($\varepsilon_s = 0.030$) is placed around the bead.
\end{problem}
\begin{solution}{P12.7}
\begin{enumerate}
  \item[(a)] At steady state, convective heat transfer from the gas to the thermocouple bead equals radiative heat loss from the small bead to the large enclosing duct walls ($F_{tc\to w} = 1$, $A_{tc}/A_w \to 0$):
  \[
    h_{\mathrm{conv}}\left(T_g - T_{tc}\right) = \varepsilon_{tc}\sigma\left(T_{tc}^4 - T_w^4\right) = (0.850)(\num{5.670374e-8})\left(450.0^4 - 350.0^4\right) = (4.8198\times 10^{-8})(\num{2.6000e10}) = \mathbf{\SI{1253.2}{\watt\per\square\meter}}
  \]
  \[
    T_g = T_{tc} + \frac{\SI{1253.2}{\watt\per\square\meter}}{\SI{65.0}{\watt\per\square\meter\per\kelvin}} = 450.0 + 19.28 = \mathbf{\SI{469.28}{\kelvin}} = \mathbf{\SI{196.13}{\celsius}}
  \]
  The unshielded thermocouple under-reads the true gas temperature by $\mathbf{-\SI{19.28}{\kelvin}}$!
  \item[(b)] Surrounding the bead with a polished radiation shield reduces the radiative loss by $>90\%$, cutting the radiation error from $-\SI{19.3}{\kelvin}$ to $<\SI{1.2}{\kelvin}$ ($T_{tc,\mathrm{shielded}} \approx \SI{468.1}{\kelvin}$).
\end{enumerate}
\end{solution}

\begin{problem}{P12.8}{Two-Surface Gray Enclosure vs.\ Blackbody Exchange Between Unequal Cylinders}{12.2}{5.c.vi}{State / National}{60 s}{Oppenheim 3-Resistor Formula}{applying parallel-plate $\frac{1}{\varepsilon_1}+\frac{1}{\varepsilon_2}-1$ without area ratio $A_1/A_2$}
A hot inner cylinder of diameter $D_1 = \SI{60.0}{\milli\meter}$ ($T_1 = \SI{360.0}{\kelvin}$, $\varepsilon_1 = 0.200$) is surrounded concentrically by a cold outer cylinder of diameter $D_2 = \SI{180.0}{\milli\meter}$ ($T_2 = \SI{300.0}{\kelvin}$, $\varepsilon_2 = 0.100$) across an evacuated gap of length $\ell = \SI{0.500}{\meter}$. Calculate (a) the exact radiative heat transfer rate $\dot{Q}_{1\to 2}$, and (b) the percentage error committed if a student uses the flat parallel-plate denominator $\frac{1}{\varepsilon_1} + \frac{1}{\varepsilon_2} - 1$ instead of weighting the outer surface resistance by $A_1/A_2 = D_1/D_2$.
\end{problem}
\begin{solution}{P12.8}
\begin{enumerate}
  \item[(a)] Inner surface area is $A_1 = \pi D_1 \ell = \pi(0.0600)(0.500) = \SI{0.094248}{\square\meter}$, and the area ratio is $A_1 / A_2 = D_1 / D_2 = 60.0 / 180.0 = 1/3$. By Lienhard Eq.~(10.27) (Eq.~\ref{eq:12:concentric-vf}):
  \[
    \frac{1}{\mathcal{F}_{1\to 2}} = \frac{1}{\varepsilon_1} + \frac{A_1}{A_2}\!\left(\frac{1}{\varepsilon_2} - 1\right) = \frac{1}{0.200} + \frac{1}{3}(10.0 - 1) = 5.00 + 3.00 = \mathbf{8.000}
  \]
  \[
    \dot{Q}_{1\to 2} = \frac{A_1 \sigma\left(T_1^4 - T_2^4\right)}{8.000} = \frac{(0.094248)(\num{5.670374e-8})\left(360.0^4 - 300.0^4\right)}{8.000} = \frac{(0.094248)(493.10)}{8.000} = \mathbf{\SI{5.809}{\watt}}
  \]
  \item[(b)] Using the flat-plate denominator $\frac{1}{0.200} + \frac{1}{0.100} - 1 = 14.00$ would predict $\dot{Q}_{\mathrm{flat}} = \frac{8.00}{14.00}(5.809) = \SI{3.320}{\watt}$ — an error of $\mathbf{-\SI{42.9}{\percent}}$!
\end{enumerate}
\end{solution}

\alsosee{\spref{sec:01.11} (\texttt{QS-11} Radiation Fast Sheet), \spref{sec:10} (Conduction Networks), \spref{sec:11} (Convection \& Evaporation), \spref{sec:13} (Transient Lumped Heating \& $T^4$ Vacuum Cooling), \spref{sec:14} (Optical Materials Table), \spref{sec:15} (Ray-Traced Ellipsoidal/CPC Concentrator).}
